\chapter{Strenge Vergleiche und Dualität}
\label{b12guide:ch:strict}

\section{Strenge und nichtstrenge Schreibweise}

\begin{definition}[Strenge Ordnung und Unvergleichbarkeit]
\label{b12guide:def:strict}
Für $x,y\in P$ sei
\[
\begin{aligned}
 x<y&\quad\Longleftrightarrow\quad x\leq y\ \land\ x\ne y,\\
 x\parallel y&\quad\Longleftrightarrow\quad
                     x\nleq y\ \land\ y\nleq x.
\end{aligned}
\]
Im zweiten Fall heißen $x$ und $y$ \emph{unvergleichbar}.
Andernfalls heißen sie \emph{vergleichbar}.
Eine \emph{strenge Halbordnung} auf einer Menge ist eine
irreflexive und transitive Relation.
\end{definition}

\begin{theorem}[Übergang zwischen beiden Schreibweisen]
\label{b12guide:thm:strict-correspondence}
Der aus $\leq$ gebildete Operator $<$ ist eine strenge
Halbordnung. Er ist asymmetrisch:
\[
 x<y\ \Rightarrow\ \neg(y<x).
\]
Umgekehrt ist für jede irreflexive und transitive Relation $\prec$
auf $P$ die durch
\[
 x\preceq y\quad\Longleftrightarrow\quad x=y\ \lor\ x\prec y
\]
definierte Relation eine Halbordnung. Beide Verfahren sind
zueinander invers: Der strenge Anteil von $\preceq$ ist $\prec$,
und das Hinzufügen der Gleichheit zu $<$ ergibt $\leq$.
\end{theorem}

\begin{theorem}[Strenges Vergleichskriterium und gemischte Transitivität]
\label{b12guide:thm:mixed}
Für $x,y,z\in P$ gilt
\[
\begin{aligned}
 x<y&\quad\Longleftrightarrow\quad x\leq y\ \land\ \neg(y\leq x),\\
 x<y\ \land\ y\leq z&\quad\Longrightarrow\quad x<z,\\
 x\leq y\ \land\ y<z&\quad\Longrightarrow\quad x<z.
\end{aligned}
\]
\end{theorem}

\begin{theorem}[Vier Vergleichsfälle]
\label{b12guide:thm:four-cases}
Für je zwei Elemente $x,y\in P$ gilt genau eine der Aussagen
\[
 x=y,\qquad x<y,\qquad y<x,\qquad x\parallel y.
\]
Unvergleichbarkeit ist symmetrisch und irreflexiv; sie muß nicht
transitiv sein. Insbesondere bedeutet $\neg(x\leq y)$ im
allgemeinen nicht $y<x$.
\end{theorem}

\begin{theorem}[Keine nichttrivialen endlichen Vergleichskreise]
\label{b12guide:thm:cycles}
Für $n\geq1$ und $x_0,\ldots,x_n\in P$ gilt
\[
 x_0\leq x_1\leq\cdots\leq x_n\leq x_0
 \quad\Longrightarrow\quad x_0=x_1=\cdots=x_n.
\]
Eine geschlossene endliche Folge solcher Vergleiche kann daher
keinen strengen Schritt enthalten. Diese Aussage setzt keine
Endlichkeit des Trägers voraus.
\end{theorem}

\section{Ordnung umkehren}

\begin{definition}[Duale Ordnung]\label{b12guide:def:dual}
Die \emph{duale Ordnung} auf demselben Träger ist definiert durch
\[
 x\leq^{\mathrm{op}}y\quad\Longleftrightarrow\quad y\leq x.
\]
Die duale halbgeordnete Menge wird mit $P^{\mathrm{op}}$ bezeichnet.
\end{definition}

\begin{theorem}[Dualität]\label{b12guide:thm:duality}
$P^{\mathrm{op}}$ ist eine halbgeordnete Menge, und
$(P^{\mathrm{op}})^{\mathrm{op}}=P$ einschließlich der Ordnung.
Der strenge Anteil der dualen Ordnung erfüllt
\[
 x<^{\mathrm{op}}y\quad\Longleftrightarrow\quad y<x.
\]
Vergleichbarkeit und Unvergleichbarkeit bleiben beim Dualisieren
unverändert. Jede für alle Halbordnungen gültige Aussage bleibt
nach durchgängiger Umkehrung der Ordnung gültig.
\end{theorem}
